\chapter{Conclusion}
\label{ch:conclusion-battery}

The central lesson of this battery thesis is straightforward: when telemetry
degrades, a controller can satisfy its observed constraints while violating
the real battery envelope.  That hidden gap is measurable, operationally
meaningful, and correctable.

\section{Battery Outcome}

The current ORIUS battery package now combines:

\begin{itemize}
  \item a battery-domain formulation of the hidden safety problem,
  \item a DC3S runtime that turns telemetry quality into conservative action
    adjustment,
  \item a hardened artifact surface across benchmark, latency, calibration,
    aging, blackout, fleet, probing, and HIL layers, and
  \item a chapter-by-chapter manuscript structure aligned to those artifacts.
\end{itemize}

\section{Why The Result Matters}

For battery operation, the consequence is practical rather than rhetorical.
Unsafe dispatch is no longer something inferred only after damage appears.
It becomes something the controller can anticipate, bound, repair, and audit
step by step.  That is the contribution of ORIUS in the battery domain.

\section{Final Boundary}

This thesis does not conclude that every physical AI system is now solved.
It concludes that the battery domain provides a complete reference
instantiation of the ORIUS logic and that this reference instantiation is now
implemented as code, artifacts, and chapter-level evidence rather than as an
aspirational outline.

\section{Closing Statement}

The battery thesis therefore ends where it should: with a locked,
traceable, battery-first framework that is ready for harder validation
layers, not with an overextended universal claim.
